\section{Nyquist N-th Band Filter}
    \subsection{定義}
        \label{Nyquist N-th Band Filter の定義}
        $N\in\naturalNumbers,\;N\geq 2,\;\INyq\coloneq\cup_{k\in\integers}2\pi\parens*{\bracks{-1/N,1/N}+k}$ とする。
        係数 $h:\integers\to\complexNumbers$ の DTFT を $\HDTFT$ とする。
        但し引数は正規化角周波数である。
        次の性質を満たすフィルタを Nyquist N-th Band Filter と呼ぶ。
        \begin{enumerate}
            \item $\HDTFT(-\Omega) = \HDTFT(\Omega)$
            \item $\Omega\in\realNumbers\setminus\INyq\Rightarrow\HDTFT(\Omega)=0$
            \item $\sum_{k=0}^{N-1}\HDTFT\parens*{\Omega-\frac{k}{N}2\pi} = 1$
        \end{enumerate}
        このフィルタは $1/N$ ダウンサンプリング処理に於けるアンチエイリアシング・フィルタとして用いられる。
    \subsection{性質}
        \begin{itemize}
            \item $h$ は偶関数である
            \item $\HDTFT(\pi/N) = 1/2$
            \item $h(0)=1/N$
            \item $h(n)=0\quad(N\mid n,\;n\neq 0)$
            \item $m_1,m_2$ を 2 以上の自然数とし，$m_3=m_1m_2$ とする。$m_1$-th Band Filter の係数 $h_1$ を標本化定理に従って $m_2$ 倍に補間した後に $1/m_2$ を乗じて得られる係数 $h_3$ は $m_3$-th Band Filter の係数である。
        \end{itemize}
        以下で各性質を導出する。
        \subsubsection{性質1}
            \begin{align*}
                h(-n) &= \frac{1}{2\pi}\integrate{-\pi}{\pi}{\HDTFT(\Omega)\exp\parens*{-i\Omega n}}{}{\Omega} = \frac{1}{2\pi}\integrate{-\pi}{\pi}{\HDTFT(-\Omega)\exp\parens*{i(-\Omega) n}}{}{\Omega} \\
                &= -\frac{1}{2\pi}\integrate{\pi}{-\pi}{\HDTFT(\Omega')\exp\parens*{i\Omega' n}}{}{\Omega'} \quad (\text{変数変換：}\Omega = -\Omega') \\
                &= \frac{1}{2\pi}\integrate{-\pi}{\pi}{\HDTFT(\Omega')\exp\parens*{i\Omega' n}}{}{\Omega'} = h(n)
            \end{align*}
        \subsubsection{性質2}
            定義の条件 3 で $\Omega=\pi/N$ として定義の条件 2 を適用すると $1 = \HDTFT(\pi/N) + \HDTFT(-\pi/N)$ である。
            これと定義の条件 2 より性質が従う。
        \subsubsection{性質3,4}
            定義の条件 3 の両辺を逆 DTFT する。任意の整数 $n$ に対して次式が成り立つ。
            \begin{align*}
                \text{RHS} &= \frac{1}{2\pi}\integrate{-\pi}{\pi}{\exp(i\Omega n)}{}{\Omega} \\
                \text{LHS} &= \sum_{k=0}^{N-1} \IDTFT{\Omega\mapsto\HDTFT\parens*{\Omega-\frac{k}{N}2\pi}}(n) = \sum_{k=0}^{N-1} \IDTFT{\DTFT{m\mapsto h(m)\exp\parens*{i\frac{k}{N}2\pi m}}}(n) \\
                &= \sum_{k=0}^{N-1}h(n)\exp\parens*{i\frac{k}{N}2\pi n} = h(n)\sum_{k=0}^{N-1}\exp\parens*{i\frac{k}{N}2\pi n}
            \end{align*}
            上式に於いて $n=0$ のときは $1 = Nh(0)$ であるので，$h(0)=1/N$ が従う。
            また，$n=mN\;(m\in\integers,\;m\neq 0)$ のとき $0 = Nh(n)$ であるので，$h(n)=0$ が従う。
        \subsubsection{性質5}
            $h_1$ の DTFT を $\HDTFTwithSs{1}$ とする。
            $h_1$ を標本化定理に従って連続時間信号に拡張したものを $\hc$ とし，その Fourier 変換を $\Hc$ とする（引数は正規化角周波数）。
            \ref{DTFTとエイリアシングとの関係} より，$\Hc$ は $[-\pi,\pi]$ に於いて $(\Ts/\sqrt{2\pi})\HDTFTwithSs{1}$ に等しく，$(-\infty,\pi),(\pi,\infty)$ に於いて 0 である。
            $h_3$ は $\hcWithSs{2}:t\mapsto\hc(t/m_2)/m_2$ を標本化したものと等しい。
            $\HcWithSs{2}\coloneq\FT{\hcWithSs{2}}$ とすると，$\HcWithSs{2}(\Omega) = \Hc(m_2\Omega)$ が成り立つ。
            再び \ref{DTFTとエイリアシングとの関係} より，$\DTFTwithArg{h_3}{\Omega} = \HDTFTwithSs{1}(m_2\Omega)$ が成り立ち，\ref{Nyquist N-th Band Filter の定義} の条件を満たす。